\chapter{Arguing Affirmative Action}

In this lecture Michael Sandel carries forward the Rawls discussion rather than beginning afresh. The question is still whether justice tracks moral desert, but now the terrain shifts from income and wealth to opportunities, hiring, and admission standards. Affirmative action matters, then, not as an isolated policy dispute, but as a fresh test of the deeper question. We begin with Cheryl Hopwood's complaint, follow the classroom arguments as they accumulate and are sorted, and only then arrive at Aristotle, whose way of tying justice to purpose gives the lecture its final shape.

\section{From Rawls to Hopwood}

Sandel opens by recalling a distinction from the previous discussion. Rawls had argued that distributive justice is not simply a matter of rewarding virtue. That opening matters because it tells us how to hear what follows: when we turn to admissions, we are not merely asking who has the better file; we are asking whether admission is something one morally deserves.

\begin{equation}
\text{moral desert} \neq \text{entitlements to legitimate expectations}
\end{equation}

\begin{equation}
\text{distributive justice}
\Longrightarrow
\text{opportunities, hiring decisions, admission standards}
\end{equation}

Now Sandel gives us a case before he gives us a theory. Cheryl Hopwood worked her way through school, did not come from an affluent family, put herself through community college and California State University at Sacramento, moved to Texas, became a resident, took the law school admissions test, and applied to the University of Texas Law School. She had, as Sandel emphasizes, a serious academic record:

\begin{align}
\text{Hopwood GPA} &\approx 3.8,\\
\text{Texas African-American + Mexican-American share} &\approx 40\%.
\end{align}

The law school, meanwhile, defended a race-conscious admissions policy by appealing to diversity. The transcript gives only a schematic account of the admissions metric, so we state it cautiously:

\begin{equation}
\text{academic index} \sim (\text{grades},\ \text{test scores})
\end{equation}

What generated Hopwood's complaint was not merely that race was considered, but that the consideration changed outcomes.

\begin{equation}
\text{lower academic index admitted}
\qquad
\text{while Hopwood rejected}
\end{equation}

Here Sandel makes one of the lecture's essential pivots. We are asked to put aside the legal question and judge the case morally. Is the policy fair? Does Hopwood have a legitimate complaint? The lecture does not start by telling us the answer. It begins by dividing the room.

\subsection{Question \& Answer}

Does Cheryl Hopwood have a legitimate complaint if applicants with lower academic indices were admitted under a race-conscious policy?

She plainly does if we assume that admission ought to track grades and test scores, or at least the academic index built from them. Under that assumption, the policy looks like a penalty imposed because she is white. But Sandel uses the case to expose the assumption rather than endorse it. The deeper question is already in the air: who says that the academic index alone defines what admission is for?

\section{Arbitrariness, Unequal Preparation, and Measured Potential}

The first strong defense of Hopwood comes from Bree. Race, she says, is an arbitrary factor. Hopwood cannot control being white. That makes race unlike a score or grade one has worked to earn. The principle behind the objection is clear: admissions should not reward or penalize features outside the applicant's control.

Sandel's reply is not to deny the force of the principle, but to look more carefully at what grades and scores actually measure. Anisha points to unequal schooling. If educational opportunity is unequal, then a test score can fail to represent the ability or promise of the student who received it. That point can be rendered only schematically, but the scheme is faithful to the lecture:

\begin{equation}
\text{measured academic score} \neq \text{true academic potential}
\qquad
\text{under unequal educational opportunity}
\end{equation}

This is already a more subtle argument than a simple appeal to preference. It says that admissions may still aim at academic promise and scholarly potential; they may simply need a better reading of what the numbers mean.

\paragraph{Worked derivation.} Sandel's reformulation of Anisha's point can be written as a short chain:
\begin{enumerate}
\item The admissions office wants to estimate academic promise.
\item Grades and test scores are used as proxies for that promise.
\item Unequal educational opportunity affects the meaning of those proxies.
\item Therefore a score cannot always be read at face value.
\item A just admissions policy may correct for unequal preparation while still aiming at academic promise.
\end{enumerate}

Sandel then sharpens the issue by changing the facts. Suppose two candidates have equally strong grades and test scores and both went to first-rate schools. If race still counts there, then we have moved beyond the corrective argument. We are no longer only reinterpreting the numbers; we are appealing to another principle.

\subsection{Question \& Answer}

Can affirmative action be justified simply as a correction for the distorted meaning of grades and test scores?

Yes, but only in the limited sense Sandel draws out here. On this argument, academic promise remains the standard; the correction lies in how we interpret the evidence. This is why the argument is narrower than it first appears. The moment race is counted even where unequal preparation has been set aside, we are no longer merely correcting the measure. We are asking what else a university is entitled to seek.

\section{Compensation, Legacy, and the Diversity Turn}

Once that narrower route is on the table, the discussion broadens. David offers a second defense: affirmative action as compensation for past injustice, especially slavery and segregation. Kate objects that this tries to fix results rather than causes. If the real problem lies in unequal education and upbringing, then justice should repair those upstream conditions rather than rework admissions decisions at the end.

The discussion then widens again. Hannah compares race-conscious admissions with nepotism and legacy admissions. If we tolerate advantages that come from family background, social networks, and parental attendance, it becomes harder to claim that race is the uniquely improper departure from merit. Ted replies that race-based decisions are inherently unfair. Da answers with the strongest rejoinder available in the room: a great many advantages already rest on family background, schooling, and inherited circumstances over which no one has control.

This part of the lecture matters because Sandel does not rush to classification. He allows the arguments to accumulate. Reparations, educational disadvantage, legacy admissions, social background, diversity in the classroom: all of these enter before the map is drawn. That gives the later classification its point.

\subsection{Question \& Answer}

Is affirmative action compensation for past injustice, and if so is that fair to present applicants?

The compensatory argument has real moral force. Past wrongs do not stay in the past; they shape present opportunity. But Sandel immediately identifies the hardest objection. Is it fair to ask Hopwood to bear a cost for injustices in which she was not personally implicated? Once the question is put that way, the issue becomes one of group responsibility and responsibility across time. Sandel marks that problem clearly, then brackets it. He does not deny its importance; he sets it aside so that the lecture can turn to the third argument.

\section{Three Arguments, Diversity, and the Return of Rights}

At this point Sandel deliberately stops the back-and-forth and reorganizes the terrain. If we have listened carefully, he says, three arguments have emerged in defense of counting race and ethnicity in admissions.

\begin{equation}
\text{affirmative action defenses}
=
\{\text{corrective},\ \text{compensatory},\ \text{diversity}\}
\end{equation}

The first is Anisha's argument: correcting for the effects of unequal preparation. The important feature of this argument is that it remains consistent with the view that academic promise and scholarly potential should be the governing standard. The office simply needs more than raw scores and grades to estimate that promise well.

The second is the compensatory argument: affirmative action as a response to historic wrongs. Here the appeal is backward-looking. Admissions become one site among others where society might attempt to right old injustices.

The third is the diversity argument, and this one changes the structure of the debate. It does not ask only how to read the file, and it does not ask only how to compensate for the past. It asks what purpose the university serves.

\begin{equation}
\text{diversity argument}
=
\text{educational experience} + \text{wider civic mission}
\end{equation}

Sandel is explicit that the diversity argument has two aspects. One concerns the educational experience of the students themselves: a diverse student body changes what everyone learns. The other concerns the wider society. In the Hopwood case, the University of Texas Law School argued that its task was to train lawyers, judges, legislators, and public officials who would reflect the backgrounds and experiences of the state it served.

This is why Sandel turns to the Bakke case and to Harvard's brief. Harvard's claim was not that scholarly excellence had ceased to matter, but that it had never been the sole criterion of admission. Geography, talents, rural and urban background, artistic and athletic gifts, all had long counted as sources of diversity. Race, in that argument, became one more ``plus'' factor in a class already assembled by more than one measure.

Now the lecture reaches its deepest Rawlsian difficulty. The diversity argument is an argument in the name of social purpose and the common good. But Rawls does not believe that common-good arguments automatically override individual rights.

\begin{equation}
\text{common good} \not\Longrightarrow \text{automatic override of individual rights}
\end{equation}

So what right might Hopwood claim? Perhaps the right to be considered according to factors within her control. Perhaps the right to be judged by achievements, efforts, and excellences rather than by a social purpose external to her.

Sandel then brings us back, with some force, to the opening distinction. The answer from the other side is that nobody morally deserves admission according to a mission-neutral standard, because there is no such standard prior to the institution's self-definition.

\begin{equation}
\text{no one morally deserves admission}
\end{equation}

\begin{equation}
\text{institutional mission}
\Longrightarrow
\text{admissions criteria}
\Longrightarrow
\text{legitimate expectations}
\end{equation}

Once the institution defines its purpose and its criteria, those who satisfy them may have legitimate expectations. But no one morally deserves that the institution define its mission in a way that happens to prize the qualities he or she already possesses in abundance.

\subsection{Question \& Answer}

What exactly are the three distinct arguments for affirmative action, and how does the diversity argument differ from the other two?

The corrective argument keeps academic promise as the governing aim and adjusts the interpretation of evidence. The compensatory argument justifies preference by reference to past wrongs. The diversity argument, by contrast, is telegraphed already as an argument about institutional purpose. That is what makes it the most ambitious of the three. It is also what makes it the most vulnerable to a rights-based objection.

\subsection{Question \& Answer}

Does an applicant have a right to be judged only by achievements and factors within her control?

Sandel lets the force of that complaint stand. But he also exposes its instability. Family background, schooling, inherited opportunities, and even many of the conditions under which achievement becomes possible are themselves beyond our control. So the principle cannot be applied cleanly. The lecture does not dissolve the complaint; it shows why the complaint pushes us back toward the deeper issue of desert and entitlement.

\section{Reset, Mission, and the Historical Challenge}

The lecture then passes through a visible seam. We hear a brief anticipatory line about distributing flutes, and only after that does the next session open with a recap: when we ended last time, three arguments had emerged. The reset matters. It does not merely repeat the discussion; it sharpens the unresolved issue.

The recap first restates the three arguments. Then Sandel returns to the diversity argument in its civic form. The University of Texas Law School, he says, defended the policy by claiming that its social mission was to produce leaders in law and politics who would reflect the racial and ethnic composition of Texas. This is precisely the point at which the Rawlsian challenge re-enters: can an institution deny admission to an individual for the sake of a wider mission?

Da now presses on Hannah's earlier claim that a private institution may define its mission however it wants. Sandel's response is a classic pressure test. What about the University of Texas Law School in the 1950s, when it defended segregation by reference to its mission? What about Harvard in the 1930s, when anti-Jewish quotas were defended by reference to the kinds of professions Harvard aimed to supply?

The question is no longer whether institutions have missions. Of course they do. The question is whether any mission an institution chooses can be morally authoritative.

Hannah's reply is that present-day affirmative action is inclusionary whereas the older policies were exclusionary. Sandel then refines the challenge again. Perhaps the difference lies in malice: the older policies expressed a judgment that African Americans or Jews were less worthy, whereas present-day affirmative action does not. But if that is the answer, then the lecture has identified something morally important without yet settling the deeper question. For even without malice, applicants may still be treated instrumentally, as useful to a social purpose rather than honored for merit or desert.

This is where Sandel introduces the mock admissions letters. The unsuccessful letter would say, in effect: you are not at fault, but society does not presently need the qualities you offer. The successful letter would say: lucky for you, your traits happen to be useful to society just now, so we propose to exploit them for that purpose. The joke is pointed. It reveals a moral oddness in the view that admissions is wholly detached from desert, even if entitlement follows once the rules are fixed.

Whatever the exact turns in the noisy stretch of transcript that follows, the stable philosophical point is clear. The impulse to detach justice from moral desert is not confined to Rawls. Rights-oriented theories, egalitarian and libertarian alike, tend to agree that justice is not centrally about honoring virtue. Sandel turns to Aristotle precisely to test that shared assumption.

\subsection{Question \& Answer}

Is there a principled difference between today's diversity rationale and earlier exclusionary appeals to institutional mission?

The lecture explores two candidate differences: inclusion rather than exclusion, and the absence of malicious or demeaning judgment. Both are morally significant. But Sandel's challenge is that neither by itself fully answers the question. Once admissions are structured by institutional mission, we still need a way of judging the mission itself. At that point the language of rights and entitlements begins to seem incomplete, and the lecture is ready for Aristotle.

\section{Aristotle, Flutes, and Teleological Justice}

Sandel now changes the framework. To question the shared modern assumption, he turns to a thinker who explicitly ties justice to honor, merit, and moral desert. Aristotle's view is powerful because it begins from an intuition many of us still feel, even after modern theory has tried to talk us out of it.

\begin{equation}
\text{justice} = \text{giving each person his or her due}
\end{equation}

But that formula is only the beginning. What is a person's due? Aristotle's answer is that justice always involves a relation between the thing distributed and the persons among whom it is distributed.

\begin{equation}
\text{justice involves two factors: things and persons}
\end{equation}

\begin{equation}
\text{persons equal in the relevant respect}
\Longrightarrow
\text{equal things assigned}
\end{equation}

The hard question, then, is the one Sandel immediately asks: equal in what respect? Aristotle's answer is that the relevant respect depends on the sort of thing at issue. Justice is therefore not independent of purpose. It depends on the point of the practice and the excellence appropriate to it.

\begin{equation}
\text{telos of the practice}
\Longrightarrow
\text{criterion of just discrimination}
\end{equation}

That is why the flute example is so central. If the thing being distributed is a set of flutes, the relevant question is not who is richest, noblest, most beautiful, or most influential. The relevant question is who can use a flute as a flute ought to be used.

\begin{equation}
\text{flutes}
\Longrightarrow
\text{best flute players get the best flutes}
\end{equation}

\paragraph{Worked example: the allocation of flutes.} Sandel's reconstruction of Aristotle can be stated step by step:
\begin{enumerate}
\item Identify the good being distributed: flutes.
\item Ask what flutes are for: they are for being played well.
\item From the purpose of the good, identify the relevant excellence: flute playing.
\item Allocate the best flutes to those with the relevant excellence.
\end{enumerate}

The structure can be written more compactly as

\begin{align}
\text{thing distributed} &\Longrightarrow \text{relevant excellence},\\
\text{relevant excellence} &\Longrightarrow \text{just recipient}.
\end{align}

And that immediately excludes rival criteria:

\begin{equation}
\text{wealth, birth, beauty, chance}
\not\Longrightarrow
\text{relevant basis for flute allocation}
\end{equation}

Sandel pauses here to block the most natural modern interpretation. Aristotle is not making a utilitarian point. The reason the best flutes should go to the best flute players is not merely that everyone will enjoy the resulting music more, though that may happen. The reason is that this allocation fits the purpose of the good itself.

\subsection{Question \& Answer}

Why should the best flutes go to the best flute players, and why is that answer not utilitarian?

A utilitarian answer would appeal to the consequences: better music, more pleasure, more overall benefit. Aristotle's answer appeals instead to the nature and end of the practice. Flutes are for being played well. So the just allocation is the one that honors the excellence appropriate to flute playing. The good consequences may follow, but they are not the ground of the judgment.

This is where Sandel introduces the language of \emph{telos}.

\begin{equation}
\text{teleological reasoning}
=
\text{reasoning from the goal or end}
\end{equation}

The point may seem strange only because modern thought has grown suspicious of teleological explanation. Sandel's tennis-court example restores the intuition. Who should have priority for the best tennis courts? Not the richest, not the most powerful, not even the greatest scientist on the faculty, if scientific greatness is irrelevant to tennis. The courts are for excellent play, so the relevant excellence is excellence in tennis.

He then widens the frame one final time. In Aristotle's world, not only social practices but nature itself was understood teleologically. Sandel's brief Winnie-the-Pooh example is meant to show that this mode of thought still has a natural grip on us: buzzing points to bees, bees to honey, honey to what it is for. Whether or not we accept teleology in nature, the moral question remains alive in social life.

That is why the lecture closes by returning, not away from affirmative action, but back through it.

\begin{equation}
\text{affirmative action dispute}
\Longrightarrow
\text{disagreement about the telos of university education}
\end{equation}

If we ask what a university is for, what excellences it ought to cultivate, and what sort of community it ought to form, then we have not changed the subject. We have reached the deepest level of the subject. Aristotle enters the lecture not as an appendage, but as the thinker who makes plain what the argument about affirmative action had really been about all along.

\section{Summary}

The lecture begins with Rawls and ends with Aristotle, but the path is continuous. Sandel first reopens the distinction between moral desert and legitimate expectations, then uses Hopwood's complaint to test that distinction in the setting of admissions. Out of the classroom exchange three defenses of affirmative action emerge: corrective, compensatory, and diversity. The diversity argument proves the most philosophically ambitious, because it appeals to the mission of the institution and so raises the Rawlsian question whether common-good reasoning may override individual claims.

The historical challenge to mission-based admissions, together with the mock admissions letters, then makes the moral oddness of a purely anti-desert view impossible to ignore. That pressure opens the way to Aristotle, whose teleological account of justice ties allocation to purpose and relevant excellence. The final implication is plain. The dispute over affirmative action is, at bottom, a dispute over what a university is for. Justice here cannot be understood without first asking about the end of the practice itself.